\documentclass[16pt, letterpaper]{article}
\usepackage[T2A]{fontenc}
\usepackage{amsfonts}
\usepackage{amssymb}

\title{Обязательные формулировки}
\author{Курбанов Сергей}
\date{Математический анализ, 3 семестр}

\begin{document}

\maketitle 

\section*{Равномерная сходимость}
\textbf{Определение.} Пусть $X$ -- множество, $f_n$, $f : X \rightarrow \mathbb{R}$ (или $\mathbb{C}$). Говорят, что последовательность $\{f_n\}$ \textit{сходится} к функции $f$ на множестве $X$ \textit{равномерно}, и пишут $$f_n \rightrightarrows\limits_{n \rightarrow \infty} f (X),$$ если $$\forall \epsilon > 0 \quad \exists N \in \mathbb{N} \quad \forall n > N \quad \forall x \in X \quad |f_n(x) - f(x)| < \epsilon$$

\section*{Равномерная норма}
\textbf{Определение.} Пусть $X$ -- множество, $f: X \to \mathbb{R}$ (или $\mathbb{C}$). Величина $$\|f\| = sup\limits_{x \in X} |f(x)|$$ называется \textit{равномерной} или \textit{чебышевской нормой} функции $f$.

\end{document}

